% GENERATED FILE. Edit canonical lessons, then run npm run generate:books.
% canonical-source-hash: fnv1a64:e1419d276255edc0
% canonical-lessons: JA-C195-kita, JA-C195-tsukiatari, JA-C195-mukai, JA-C195-sugiru, JA-C195-modoru, JA-R195-review-190-195

\chapter{North, Opposite and Past}
\label{ch:ja-north-opposite-and-past}
\hlchaptermodality{195}

\begin{chapteropening}
\textbf{By the end of this chapter:} \emph{I can say where a place is from a landmark.}

Complete the last lesson of chapter 195: I can say where a place is from a landmark.
\end{chapteropening}

\section[\texorpdfstring{\ja{きた} (kita)}{きた (kita)}]{\ja{きた} (kita) — north}

\label{lesson:JA-C195-kita}



\begin{quote}
Before the new one: say the Japanese for to pass through, then the Japanese for museum.
\end{quote}

\subsection*{You'll want to know: \ja{きた}}
\textbf{\ja{きた}} — \emph{kita} — \textquotedblleft{}north\textquotedblright{}.

\begin{quote}
\textbf{\ja{えきの} \ja{きたに} \ja{こうえんが} \ja{あります。}} — \emph{eki no kita ni kōen ga arimasu} — \textquotedblleft{}There is a park north of the station.\textquotedblright{}
\end{quote}

\begin{cousinweb}
\textbf{\ja{きた}} (\emph{kita}) is a native Japanese word (wago), inherited rather than borrowed from Chinese; its kanji was attached later as a kun'yomi (JMdict, KANJIDIC2).
\end{cousinweb}

\begin{grammarlens}[title={What you've built}]
One more word for everyday life at A2.
\end{grammarlens}

\subsection*{Your turn}
\begin{itemize}
\raggedright
  \item \emph{Say it:} \emph{kita}
  \item \emph{Say it:} \emph{kita}, once more
  \item \emph{Say it:} \emph{eki no kita ni kōen ga arimasu} — There is a park north of the station.
  \item \emph{Recall:} say the Japanese for to pass through, then the Japanese for museum, then say \emph{kita} again
\end{itemize}

\begin{tcolorbox}[breakable,colback=teal!4,colframe=teal!35!black,title={Before you move on}]
What does \ja{きた} mean? (\textquotedblleft{}North\textquotedblright{}.) Say it once more.
\end{tcolorbox}
\section[\texorpdfstring{\ja{つきあたり} (tsukiatari)}{つきあたり (tsukiatari)}]{\ja{つきあたり} (tsukiatari) — end of the street; dead end}

\label{lesson:JA-C195-tsukiatari}



\begin{quote}
Before the new one: say the Japanese for museum, then the Japanese for north.
\end{quote}

\subsection*{You'll want to know: \ja{つきあたり}}
\textbf{\ja{つきあたり}} — \emph{tsukiatari} — \textquotedblleft{}end of the street; dead end\textquotedblright{}.

\begin{quote}
\textbf{\ja{この} \ja{みちの} \ja{つきあたりに} \ja{ゆうびんきょくが} \ja{あります。}} — \emph{kono michi no tsukiatari ni yūbinkyoku ga arimasu} — \textquotedblleft{}The post office is at the end of this road.\textquotedblright{}
\end{quote}

\begin{cousinweb}
\textbf{\ja{つきあたり}} (\emph{tsukiatari}) is native Japanese, built from \emph{tsuki}, \textquotedblleft{}thrusting\textquotedblright{}, and \emph{atari}, \textquotedblleft{}hitting\textquotedblright{}, from \emph{ataru}, \textquotedblleft{}to hit\textquotedblright{}: the place where a street runs into something (JMdict, KANJIDIC2).
\end{cousinweb}

\begin{grammarlens}[title={What you've built}]
One more word for everyday life at A2.
\end{grammarlens}

\subsection*{Your turn}
\begin{itemize}
\raggedright
  \item \emph{Say it:} \emph{tsukiatari}
  \item \emph{Say it:} \emph{tsukiatari}, once more
  \item \emph{Say it:} \emph{kono michi no tsukiatari ni yūbinkyoku ga arimasu} — The post office is at the end of this road.
  \item \emph{Recall:} say the Japanese for museum, then the Japanese for north, then say \emph{tsukiatari} again
\end{itemize}

\begin{tcolorbox}[breakable,colback=teal!4,colframe=teal!35!black,title={Before you move on}]
What does \ja{つきあたり} mean? (\textquotedblleft{}End of the street; dead end\textquotedblright{}.) Say it once more.
\end{tcolorbox}
\section[\texorpdfstring{\ja{むかい} (mukai)}{むかい (mukai)}]{\ja{むかい} (mukai) — across the street; the place opposite}

\label{lesson:JA-C195-mukai}



\begin{quote}
Before the new one: say the Japanese for north, then the Japanese for end of the street.
\end{quote}

\subsection*{You'll want to know: \ja{むかい}}
\textbf{\ja{むかい}} — \emph{mukai} — \textquotedblleft{}across the street; the place opposite\textquotedblright{}.

\begin{quote}
\textbf{\ja{ぎんこうの} \ja{むかいに} \ja{やっきょくが} \ja{あります。}} — \emph{ginkō no mukai ni yakkyoku ga arimasu} — \textquotedblleft{}There is a pharmacy across from the bank.\textquotedblright{}
\end{quote}

\begin{cousinweb}
\textbf{\ja{むかい}} (\emph{mukai}) is native Japanese, inherited rather than borrowed, derived from \emph{mukau}, \textquotedblleft{}to face\textquotedblright{} (JMdict, KANJIDIC2).
\end{cousinweb}

\begin{grammarlens}[title={What you've built}]
One more word for everyday life at A2.
\end{grammarlens}

\subsection*{Your turn}
\begin{itemize}
\raggedright
  \item \emph{Say it:} \emph{mukai}
  \item \emph{Say it:} \emph{mukai}, once more
  \item \emph{Say it:} \emph{ginkō no mukai ni yakkyoku ga arimasu} — There is a pharmacy across from the bank.
  \item \emph{Recall:} say the Japanese for north, then the Japanese for end of the street, then say \emph{mukai} again
\end{itemize}

\begin{tcolorbox}[breakable,colback=teal!4,colframe=teal!35!black,title={Before you move on}]
What does \ja{むかい} mean? (\textquotedblleft{}Across the street; the place opposite\textquotedblright{}.) Say it once more.
\end{tcolorbox}
\section[\texorpdfstring{\ja{すぎる} (sugiru)}{すぎる (sugiru)}]{\ja{すぎる} (sugiru) — to go past; to pass, to exceed}

\label{lesson:JA-C195-sugiru}



\begin{quote}
Before the new one: say the Japanese for end of the street, then the Japanese for across the street.
\end{quote}

\subsection*{You'll want to know: \ja{すぎる}}
\textbf{\ja{すぎる}} — \emph{sugiru} — \textquotedblleft{}to go past; to pass, to exceed\textquotedblright{}.

\begin{quote}
\textbf{\ja{あの} \ja{しんごうを} \ja{すぎます。}} — \emph{ano shingō o sugimasu} — \textquotedblleft{}We go past that traffic light.\textquotedblright{}
\end{quote}

Polite \textbf{\ja{すぎます}}, \emph{sugimasu}; an ichidan verb.

\begin{cousinweb}
\textbf{\ja{すぎる}} (\emph{sugiru}) is a native Japanese word (wago), inherited rather than borrowed from Chinese; its kanji was attached later as a kun'yomi; JMdict marks at least one of its senses as usually written in kana (JMdict, KANJIDIC2).
\end{cousinweb}

\begin{grammarlens}[title={What you've built}]
One more word for everyday life at A2.
\end{grammarlens}

\subsection*{Your turn}
\begin{itemize}
\raggedright
  \item \emph{Say it:} \emph{sugiru}
  \item \emph{Say it:} \emph{sugiru}, once more
  \item \emph{Say it:} \emph{ano shingō o sugimasu} — We go past that traffic light.
  \item \emph{Recall:} say the Japanese for end of the street, then the Japanese for across the street, then say \emph{sugiru} again
\end{itemize}

\begin{tcolorbox}[breakable,colback=teal!4,colframe=teal!35!black,title={Before you move on}]
What does \ja{すぎる} mean? (\textquotedblleft{}To go past; to pass, to exceed\textquotedblright{}.) Say it once more.
\end{tcolorbox}
\section[\texorpdfstring{\ja{もどる} (modoru)}{もどる (modoru)}]{\ja{もどる} (modoru) — to go back; to return}

\label{lesson:JA-C195-modoru}



\begin{quote}
Before the new one: say the Japanese for across the street, then the Japanese for to go past.
\end{quote}

\subsection*{You'll want to know: \ja{もどる}}
\textbf{\ja{もどる}} — \emph{modoru} — \textquotedblleft{}to go back; to return\textquotedblright{}.

\begin{quote}
\textbf{\ja{ちょっと} \ja{えきに} \ja{もどります。}} — \emph{chotto eki ni modorimasu} — \textquotedblleft{}I am going back to the station for a moment.\textquotedblright{}
\end{quote}

Polite \textbf{\ja{もどります}}, \emph{modorimasu}; a godan verb. Its transitive partner \textbf{\ja{もどす}} is taught already.

\begin{cousinweb}
\textbf{\ja{もどる}} (\emph{modoru}) is a native Japanese word (wago), inherited rather than borrowed from Chinese; its kanji was attached later as a kun'yomi, and it shares its stem with \emph{modosu}, \textquotedblleft{}to put back\textquotedblright{} (JMdict, KANJIDIC2).
\end{cousinweb}

\begin{grammarlens}[title={What you've built}]
One more word for everyday life at A2.
\end{grammarlens}

\subsection*{Your turn}
\begin{itemize}
\raggedright
  \item \emph{Say it:} \emph{modoru}
  \item \emph{Say it:} \emph{modoru}, once more
  \item \emph{Say it:} \emph{chotto eki ni modorimasu} — I am going back to the station for a moment.
  \item \emph{Recall:} say the Japanese for across the street, then the Japanese for to go past, then say \emph{modoru} again
\end{itemize}

\begin{tcolorbox}[breakable,colback=teal!4,colframe=teal!35!black,title={Before you move on}]
What does \ja{もどる} mean? (\textquotedblleft{}To go back; to return\textquotedblright{}.) Say it once more.
\end{tcolorbox}
\section[(dialogue)]{Words from chapters 190–195 — a review}

\label{lesson:JA-R195-review-190-195}



\begin{quote}
Say the Japanese for building and shopping street.
\end{quote}

\subsection*{Your turn}
Say these aloud:

\begin{itemize}
\raggedright
  \item building, shopping street, department store, to enter, to be open
  \item to close, market, liquor store, pharmacy, to find
  \item police box, police officer, police, to ask, guidance
  \item shrine, temple, cinema, art museum, to be visible
  \item right, straight ahead, to turn, to pass through, museum
  \item north, end of the street, across the street, to go past, to go back
\end{itemize}

\begin{tcolorbox}[breakable,colback=teal!4,colframe=teal!35!black,title={Before you move on}]
Building? (\textbf{\ja{たてもの}}, \emph{tatemono}.) Shopping street? (\textbf{\ja{しょうてんがい}}, \emph{shōtengai}.) Department store? (\textbf{\ja{ひゃっかてん}}, \emph{hyakkaten}.) To enter? (\textbf{\ja{はいる}}, \emph{hairu}.) To be open? (\textbf{\ja{あく}}, \emph{aku}.) To close? (\textbf{\ja{しまる}}, \emph{shimaru}.) Market? (\textbf{\ja{いちば}}, \emph{ichiba}.) Liquor store? (\textbf{\ja{さかや}}, \emph{sakaya}.) Pharmacy? (\textbf{\ja{やっきょく}}, \emph{yakkyoku}.) To find? (\textbf{\ja{みつける}}, \emph{mitsukeru}.) Police box? (\textbf{\ja{こうばん}}, \emph{kōban}.) Police officer? (\textbf{\ja{おまわりさん}}, \emph{omawarisan}.) Police? (\textbf{\ja{けいさつ}}, \emph{keisatsu}.) To ask? (\textbf{\ja{たずねる}}, \emph{tazuneru}.) Guidance? (\textbf{\ja{あんない}}, \emph{annai}.) Shrine? (\textbf{\ja{じんじゃ}}, \emph{jinja}.) Temple? (\textbf{\ja{てら}}, \emph{tera}.) Cinema? (\textbf{\ja{えいがかん}}, \emph{eigakan}.) Art museum? (\textbf{\ja{びじゅつかん}}, \emph{bijutsukan}.) To be visible? (\textbf{\ja{みえる}}, \emph{mieru}.) Right? (\textbf{\ja{みぎ}}, \emph{migi}.) Straight ahead? (\textbf{\ja{まっすぐ}}, \emph{massugu}.) To turn? (\textbf{\ja{まがる}}, \emph{magaru}.) To pass through? (\textbf{\ja{とおる}}, \emph{tōru}.) Museum? (\textbf{\ja{はくぶつかん}}, \emph{hakubutsukan}.) North? (\textbf{\ja{きた}}, \emph{kita}.) End of the street? (\textbf{\ja{つきあたり}}, \emph{tsukiatari}.) Across the street? (\textbf{\ja{むかい}}, \emph{mukai}.) To go past? (\textbf{\ja{すぎる}}, \emph{sugiru}.) To go back? (\textbf{\ja{もどる}}, \emph{modoru}.)
\end{tcolorbox}
